%!TEX root = ../main.tex
%----------------------------------------------------------------------
%  Section 4: Estimation and Implementation
%  Paper:   Theory-Informed Kernel Ridge Regression for Asset Pricing
%  Authors: Amedeo Andriollo and Juan F. Imbet
%----------------------------------------------------------------------

\section{Estimation and Implementation}\label{sec:estimation}

This section addresses the practical aspects of taking Theory-KRR from the mathematical framework of Sections~\ref{sec:framework}--\ref{sec:restrictions} to an implementable estimator.
We describe the kernel choice and bandwidth calibration (Section~\ref{sub:kernel_choice}), the pre-estimation of structural parameters that enter the penalty functionals (Section~\ref{sub:pre_estimation}), the cross-validation strategy for selecting the high-dimensional vector of penalty multipliers (Section~\ref{sub:cv_strategy}), computational details (Section~\ref{sub:computation}), and the rolling-window procedure for out-of-sample evaluation (Section~\ref{sub:rolling}).

%======================================================================
\subsection{Kernel Choice and Bandwidth Selection}\label{sub:kernel_choice}
%======================================================================

The kernel $k: \R^p \times \R^p \to \R$ determines the RKHS $\Hcal$ in which the estimator searches for the conditional expectation function $f^*$.
Our primary specification is the Gaussian radial basis function (RBF) kernel:
\begin{equation}\label{eq:gaussian_kernel}
k(\mathbf{x}, \mathbf{x}') = \exp\!\left(-\frac{\|\mathbf{x} - \mathbf{x}'\|^2}{2\sigma^2}\right),
\end{equation}
where $\sigma > 0$ is the bandwidth parameter and $\|\cdot\|$ denotes the Euclidean norm in $\R^p$.
The Gaussian kernel is universal: the associated RKHS is dense in the space of continuous functions on any compact subset of $\R^p$, so the approximation error $\inf_{f \in \Hcal}\|f - f^*\|$ can be made arbitrarily small.
This universality is essential for our application, where the true conditional expectation function may exhibit the nonlinearities, interactions, and threshold effects documented in the cross-section of returns.

We set the bandwidth $\sigma$ using the median heuristic:
\begin{equation}\label{eq:median_heuristic}
\sigma = \operatorname{median}\!\left\{\|\mathbf{X}_{i,t} - \mathbf{X}_{j,s}\| : (i,t) \neq (j,s)\right\}.
\end{equation}
The median heuristic places the kernel bandwidth at a scale where roughly half of the pairwise distances produce kernel evaluations above $e^{-1/2} \approx 0.61$ and half below.
This data-driven rule avoids both extremes: a bandwidth that is too small relative to the data's spread produces a nearly diagonal kernel matrix (overfitting), while a bandwidth that is too large produces a nearly constant kernel matrix (underfitting, approaching the global mean).
The median heuristic has strong theoretical support in the kernel methods literature and performs well in high-dimensional settings where manual tuning is infeasible.

Two limiting cases clarify the role of the kernel.
A linear kernel $k(\mathbf{x}, \mathbf{x}') = \mathbf{x}^\top \mathbf{x}'$ restricts $\Hcal$ to linear functions $f(\mathbf{x}) = \bm{\beta}^\top \mathbf{x}$, and Theory-KRR reduces to a penalized linear regression with structural penalties on coefficients.
This nesting (Proposition~\ref{prop:nesting}(iii)) allows us to isolate the contribution of nonlinearity from the contribution of theory-informed penalties.
At the other extreme, the Gaussian kernel with bandwidth $\sigma \to 0$ produces a kernel matrix approaching the identity, so $\hat{f}$ interpolates the data subject only to the structural penalties.
The joint configuration of $\sigma$, $\lambda_{\mathrm{stat}}$, and $\bm{\mu}$ determines the estimator's regularization, and cross-validation selects the balance that minimizes OOS prediction error.

%======================================================================
\subsection{Estimation of Structural Parameters}\label{sub:pre_estimation}
%======================================================================

Each structural penalty $C_j(f)$ depends on model-specific parameters $\theta_j$ that characterize the underlying economic model.
For consumption-based models, $\theta_j$ includes the coefficient of relative risk aversion $\gamma$, the elasticity of intertemporal substitution $\psi$, the subjective discount factor $\delta$, and the habit persistence parameter $b$.
For production-based models, $\theta_j$ governs the curvature of the adjustment cost function.
For intermediary models, $\theta_j$ parameterizes the SDF loadings on the intermediary factor.
We consider three approaches to estimating $\theta_j$, ordered from most to least principled.

\paragraph{Approach 1: Joint estimation (preferred).}
The most principled approach estimates $f$ and $\{\theta_j\}_{j=1}^J$ simultaneously:
\begin{equation}\label{eq:joint_estimation}
(\hat{f}, \hat{\theta}_1, \ldots, \hat{\theta}_J) \;=\; \argmin_{f \in \mathcal{H}, \, \theta_1, \ldots, \theta_J} \;
\frac{1}{n}\sum_s (R_s - f(\mathbf{X}_s))^2 + \lambda_{\mathrm{stat}}\|f\|_{\mathcal{H}}^2 + \sum_j \mu_j\, C_j(f; \theta_j).
\end{equation}
This optimization is convex in $f$ (given $\theta$) and generally nonconvex in $\theta$, but alternating minimization---updating $f$ given $\theta$ and then $\theta$ given $f$---converges reliably from the pre-estimation starting point.
Joint estimation is the most principled approach because it allows the structural parameters and the nonparametric function to adapt to each other: a consumption-based SDF with risk aversion $\gamma$ that is optimal in isolation may not be optimal when combined with intermediary and production penalties.

\paragraph{Approach 2: Pre-estimation (two-stage).}
For computational simplicity, one can pre-estimate $\hat{\theta}_j$ in a first stage that is separate from the KRR optimization.
For each consumption-based SDF model $j$, we estimate
\begin{equation}\label{eq:gmm_first_stage}
\hat{\theta}_j = \argmin_{\theta_j} \left\|\frac{1}{T}\sum_{t=1}^{T} \mathbf{z}_t \left[M_j(\theta_j, \mathbf{X}_t, \mathbf{X}_{t+1})\, R_{m,t+1}^{\mathrm{gross}} - 1\right]\right\|_{\hat{\mathbf{W}}^{-1}}^2,
\end{equation}
where $\mathbf{z}_t$ is a vector of instruments (lagged consumption growth, the dividend-price ratio, the term spread) and $\hat{\mathbf{W}}$ is an HAC-consistent estimate of the instrument covariance matrix.
For production-based models, we estimate adjustment cost parameters from firm-level investment regressions.
For intermediary models, the SDF loadings $(\hat{a}_j, \hat{b}_j)$ are estimated by GMM on the Euler equation $\E[M_j R^{\mathrm{gross}}] = 1$ using test assets.
The pre-estimated $\hat{\theta}_j$ is then plugged into $C_j(f)$ and treated as fixed during the second-stage KRR optimization.
This approach ignores the interaction between $f$ and $\theta$: the structural parameters estimated from aggregate moments may not be optimal for cross-sectional return prediction.

\paragraph{Approach 3: Calibration.}
For parameters with well-established values, we adopt published calibrations from the original papers.
We set $\delta = 0.99$ (monthly) in all consumption-based models and use $b = 0.87$ from \citet{campbellcochrane1999} for the external habit model.
This is the simplest approach: $C_j(f)$ is fully determined by the data and a known parameter vector, with no estimation of $\theta_j$.
The risk is that published parameter values, calibrated to match aggregate moments in a specific sample, may not be appropriate for cross-sectional return prediction.

Our baseline uses joint estimation (Approach~1), initialized from the pre-estimation values (Approach~2).
Robustness analysis compares all three approaches.
For the sensitivity analysis, we report results over a grid of values for the most influential parameters:
\begin{equation}\label{eq:sensitivity_grid}
\gamma \in \{2, 5, 10, 20\}, \qquad \delta \in \{0.95, 0.97, 0.99\}.
\end{equation}
This grid spans the range of calibrations used in the asset-pricing literature, from low-risk-aversion economies ($\gamma = 2$) to the high-risk-aversion calibrations required by rare disaster models ($\gamma = 20$).

%======================================================================
\subsection{Cross-Validation Strategy}\label{sub:cv_strategy}
%======================================================================

The full hyperparameter vector of Theory-KRR is $(\lambda_{\mathrm{stat}}, \mu_1, \ldots, \mu_J)$, where $J = 56$ in our baseline specification.
A 57-dimensional grid search is computationally infeasible.
We reduce dimensionality by exploiting the economic structure of the restrictions and adopt the three-way temporal split of \citet{gu2020empirical} (Section~1.1) for hyperparameter selection.

We group the 56 restrictions into $G = 8$ families corresponding to the eight economic model classes cataloged in Section~\ref{sec:restrictions}: consumption-based Euler equations, production-based restrictions, intermediary restrictions, information and learning restrictions, demand-system restrictions, variance and volatility restrictions, behavioral restrictions, and macro-finance restrictions.
Within each family $g$, we impose a common multiplier:
\begin{equation}\label{eq:grouped_penalties}
\mu_j = \mu_{g(j)}^{\mathrm{group}} \qquad \text{for all } j \text{ in family } g.
\end{equation}
This reduces the search space from 61 dimensions to 9: one statistical penalty $\lambda_{\mathrm{stat}}$ and eight group-level multipliers $\mu_1^{\mathrm{group}}, \ldots, \mu_8^{\mathrm{group}}$.
The grouping is motivated by the observation that restrictions within the same economic family share a common source of identification and are likely to have correlated empirical content.
The grouping is not restrictive in the dimension that matters: cross-validation can set any group's multiplier to zero, shutting down that entire family of restrictions.

Following \citet{gu2020empirical}, we adopt a rolling three-way sample split.
For each test year $\tau$ from 1987 to 2023, we define:
\begin{enumerate}
\item \textbf{Training sample} (July 1963 through year $\tau - 13$): estimate $f(\cdot)$ and all structural parameters $\hat{\theta}_j$ for each candidate hyperparameter configuration.
\item \textbf{Validation sample} (year $\tau - 12$ through year $\tau - 1$, 12 years): select the hyperparameter vector $(\hat{\lambda}_{\mathrm{stat}}, \hat{\bm{\mu}}^{\mathrm{group}})$ that minimizes average squared prediction error:
\begin{equation}\label{eq:cv_criterion}
(\hat{\lambda}_{\mathrm{stat}}, \hat{\bm{\mu}}^{\mathrm{group}}) = \argmin_{\lambda_{\mathrm{stat}}, \bm{\mu}^{\mathrm{group}}} \frac{1}{n_{\mathrm{val}}}\sum_{(i,t) \in \mathcal{S}_{\mathrm{val}}} \left(R_{i,t+1} - \hat{f}_{\mathrm{train}}(\mathbf{X}_{i,t})\right)^2,
\end{equation}
where $\hat{f}_{\mathrm{train}}$ is estimated using training data only.
\item \textbf{Testing sample} (year $\tau$, 12 months): evaluate out-of-sample performance using the selected hyperparameters.
This sample is never used for estimation or tuning.
\end{enumerate}
For the first test year $\tau = 1987$, the training sample spans 1963--1974 (12 years) and the validation sample spans 1975--1986 (12 years).
As $\tau$ advances, both the training and validation windows expand and slide forward: by 2023, the training sample spans 1963--2010 and the validation sample spans 2011--2022.
The 12-year validation window follows \citet{gu2020empirical} and provides sufficient temporal variation for reliable hyperparameter selection.
After selecting $(\hat{\lambda}_{\mathrm{stat}}, \hat{\bm{\mu}}^{\mathrm{group}})$ on the validation sample, we re-estimate $f$ on the combined training and validation data (all months through year $\tau - 1$) and predict returns for year $\tau$.
This rolling protocol yields 37 annual out-of-sample predictions spanning 1987--2023, covering the 1987 crash, the dot-com bust, the global financial crisis, and the COVID-19 pandemic.

We search the 9-dimensional hyperparameter space using coordinate descent over the group multipliers, an approach inspired by the cyclic coordinate descent algorithm that underlies the Lasso and elastic net solvers of \citet{friedman2010regularization}.
The statistical penalty $\lambda_{\mathrm{stat}}$ is first selected by standard KRR cross-validation without structural penalties ($\mu_g = 0$ for all $g$).
Holding $\lambda_{\mathrm{stat}}$ fixed, we then optimize the eight group multipliers by cycling through each group $g = 1, \ldots, G$ and performing a one-dimensional line search over a grid of candidate values:
\begin{equation}\label{eq:cd_grid}
\mu_g^{\mathrm{group}} \in \{0\} \cup \{10^{u} : u \in \mathrm{linspace}(-4, 1, 15)\}.
\end{equation}
For each candidate value of $\mu_g^{\mathrm{group}}$, all other multipliers are held at their current values, and the validation-set MSE is evaluated.
The value that yields the lowest MSE is retained, and the algorithm moves to the next group.
One full pass through all $G = 8$ groups constitutes one cycle; we repeat for up to 3 cycles and stop early if no group multiplier changes in a complete cycle.

The total number of function evaluations is at most $3 \times 8 \times 16 = 384$, compared to $k^9$ for an exhaustive grid.
Each evaluation requires solving one Nystr\"{o}m-approximated KRR system (an $m \times m$ Cholesky factorization with $m = 500$; see Appendix~\ref{app:nystrom}), followed by a single Newton-like correction step that incorporates the structural penalty gradient.
Because the 16 candidate values within each group are independent, they are evaluated in parallel across CPU cores, so each group's line search completes in the wall-clock time of a single evaluation.
On our 40-core workstation, the full coordinate descent completes in under 5 seconds per rolling window.

This approach connects to the penalized GMM literature, where sequential moment selection achieves rates comparable to the oracle that knows which moments are valid \citep{Caner2009lasso,Cheng2015select}.
Coordinate descent also has a natural economic interpretation: it asks, for each theory family in turn, ``what is the optimal strength of this family's penalty, given the penalties already imposed by other families?''
The sequential nature of the algorithm means that substitution effects between families---for instance, consumption-based and intermediary restrictions that capture similar risk exposures---are resolved automatically: the family that provides the largest marginal improvement is retained at a high multiplier, while the redundant family's multiplier is driven toward zero.

%======================================================================
\subsection{Computational Details}\label{sub:computation}
%======================================================================

By Proposition~\ref{prop:representer}, the estimator $\hat{f}(\cdot) = \sum_{s=1}^n \hat{\alpha}_s\, k(\mathbf{X}_s, \cdot)$ is parameterized by the coefficient vector $\hat{\bm{\alpha}} \in \R^n$.
The optimization problem~\eqref{eq:finite_dim} is therefore finite-dimensional: we optimize over $\bm{\alpha}$ rather than over functions in $\Hcal$.

When all active penalties are quadratic in the fitted values $\hat{\mathbf{f}} = \mathbf{K}\bm{\alpha}$, the objective is a convex quadratic in $\bm{\alpha}$, and the solution is given by the linear system~\eqref{eq:alpha_star}.
The dominant computational cost is the $O(n^3)$ matrix inversion (or, equivalently, the Cholesky decomposition of the $n \times n$ system matrix).
We solve the system using Cholesky factorization, which is numerically stable and exploits the positive definiteness guaranteed by $\lambda_{\mathrm{stat}} > 0$.

For problems involving Type~A (Euler equation) penalties, the system matrix in~\eqref{eq:foc} contains the term $\bm{\Omega}\,\mathbf{K}$, where $\bm{\Omega}$ includes the $\mathbf{A}_j^\top \mathbf{A}_j$ matrices.
These are at most rank $T$ (since $\mathbf{A}_j$ has $T$ rows), and we exploit this low-rank structure by precomputing and caching the products $\mathbf{A}_j\,\mathbf{K}$ for each active Euler-equation restriction.

When structural penalties are active, we approximate the theory-informed solution by a single Newton-like correction from the standard KRR warm start.
Specifically, let $\hat{\boldsymbol{\beta}}_{\mathrm{KRR}}$ denote the Nystr\"{o}m KRR solution (all $\mu_j = 0$), and let $\mathbf{g} = \sum_j \mu_j \, \mathbf{Z}^\top \frac{\partial C_j}{\partial \hat{\mathbf{f}}} \big|_{\hat{\mathbf{f}} = \mathbf{Z}\hat{\boldsymbol{\beta}}_{\mathrm{KRR}}}$ be the projected structural penalty gradient evaluated at the KRR fit.
The theory-informed coefficient vector is
\begin{equation}\label{eq:newton_step}
\hat{\boldsymbol{\beta}}_{\mathrm{theory}} = \hat{\boldsymbol{\beta}}_{\mathrm{KRR}} - \alpha \, (\mathbf{Z}^\top\mathbf{Z} + n\lambda_{\mathrm{stat}}\,\mathbf{I}_m)^{-1} \mathbf{g},
\end{equation}
where the step size $\alpha$ is selected from a discrete set $\{0, 0.01, 0.05, 0.1, 0.5, 1, 2\}$ by minimizing the validation-set MSE.
This one-step correction is computationally cheap---it requires one $m \times m$ Cholesky factorization plus one matrix-vector product per active restriction---and avoids the overhead of iterative L-BFGS optimization.
Since the structural penalties are small relative to the data-fit term (by construction, as the $\mu_j$ are tuned to improve validation MSE rather than to enforce exact compliance), the single Newton step provides a good approximation to the fully converged solution.

The kernel matrix $\mathbf{K}$ is $n \times n$ and must be stored in memory.
For our primary specification with $n \approx 4{,}000$ (cross-section of managed portfolios over 12 months), $\mathbf{K}$ occupies approximately 128~MB in double precision, which is manageable on standard hardware.
For the full individual-stock panel ($n > 100{,}000$), storing $\mathbf{K}$ requires over 80~GB, which exceeds typical workstation memory.

Two standard approximations address this scaling challenge.
The Nystr\"{o}m approximation selects a subset of $m \ll n$ landmark points and approximates $\mathbf{K} \approx \mathbf{K}_{nm}\,\mathbf{K}_{mm}^{-1}\,\mathbf{K}_{nm}^\top$, where $\mathbf{K}_{nm}$ is the $n \times m$ matrix of kernel evaluations between all points and the landmarks, and $\mathbf{K}_{mm}$ is the $m \times m$ kernel matrix among landmarks.
This reduces memory to $O(nm)$ and solve time to $O(nm^2)$.
Random Fourier features offer an alternative: they approximate the Gaussian kernel by an explicit feature map $\bm{\phi}: \R^p \to \R^D$ such that $k(\mathbf{x}, \mathbf{x}') \approx \bm{\phi}(\mathbf{x})^\top \bm{\phi}(\mathbf{x}')$, converting the kernel method into a linear regression in $D$-dimensional feature space.
Both approximations introduce a controllable approximation error that decreases as $m$ or $D$ increases.
We use $m = 500$ Nystr\"{o}m landmarks when the full kernel matrix is infeasible.


%======================================================================
\subsection{Rolling-Window Procedure}\label{sub:rolling}
%======================================================================

The rolling three-way split described in Section~\ref{sub:cv_strategy} defines the out-of-sample evaluation framework.
At each annual rebalancing date (December of year $\tau - 1$), the full procedure is:
(i)~estimate $\hat{f}$ on the training sample using each candidate $(\lambda_{\mathrm{stat}}, \bm{\mu}^{\mathrm{group}})$,
(ii)~select the hyperparameter vector that minimizes validation-sample MSE,
(iii)~re-estimate $\hat{f}$ on the combined training and validation data (all months through December of year $\tau - 1$), and
(iv)~predict returns for the 12 months of year $\tau$.
Portfolio weights are set at the rebalancing date and held fixed for 12 months.

The expanding training window grows from 12 years (for $\tau = 1987$) to 48 years (for $\tau = 2023$), allowing the estimator to learn from the full history of cross-sectional return patterns.
Annual rebalancing reduces computational cost (one full estimation per year rather than per month) and aligns with the holding-period assumptions of many characteristic-based strategies.

The hyperparameter vector $(\hat{\lambda}_{\mathrm{stat}}, \hat{\bm{\mu}}^{\mathrm{group}})$ is re-estimated at each rebalancing date, allowing the optimal balance between statistical and theory-informed regularization to vary over time.
Tracking $\hat{\mu}_g^{\mathrm{group}}$ across rebalancing dates reveals temporal variation in the value of each theory family---for instance, whether consumption-based restrictions become more useful during recessions or whether intermediary restrictions gain importance after financial crises.
We report the time series of estimated multipliers in Section~\ref{sec:results}.

All information used in estimation---returns, characteristics, macro state variables, and the structural parameters $\hat{\theta}_j$---is available as of the training window's end date.
Characteristics are lagged by at least one month relative to the return being predicted.
Macro variables use the vintage available at the prediction date, not subsequently revised values.
The structural parameters $\hat{\theta}_j$ are pre-estimated using only data within the training window.
